% Chapter 5: Results — Handbook 8.2.6
% ~6000 words

\section{Results}
\label{sec:results}

This chapter presents the evidence that the system works, how well it performs,
and what privacy it actually provides.  Each section ties back to the security
policy defined in Section~\ref{sec:intro-security-policy}.

\subsection{System Architecture}
\label{sec:results-architecture}

The final system had four layers: a Dioxus desktop UI (with a headless
binary for testing), an auction coordination layer (state machine, DHT
operations, MPC orchestration), the MPC tunnel proxy (bridging MP-SPDZ's
TCP to Veilid \texttt{app\_call}), and the Veilid overlay (DHT, private
routes, node identity).  There was no central server and no direct IP
connections between participants.  All persistent state was CBOR-serialised
into DHT records and all real-time communication travelled over
onion-encrypted private routes.

\subsubsection{DHT as Shared State}

The marketplace is bootstrapped from a single shared DHT record: the
\emph{master registry}.  All Lattice nodes know this record's key (derived
from the network key) and share a write keypair, so any node can add entries.
The full record hierarchy is shown in Figure~\ref{fig:dht-hierarchy}.

The DHT lifecycle (seller registration, listing creation, discovery, bidding,
and MPC route exchange) follows the record hierarchy in
Figure~\ref{fig:dht-hierarchy}.  The master registry and bid announcement
registries use a G-Set CRDT: each write merges the local set with the DHT
value via commutative union, so all nodes converge regardless of write order.

\subsection{Auction Protocol}
\label{sec:results-protocol}

The auction lifecycle, following the state machine in
Figure~\ref{fig:auction-fsm}, proceeds in seven phases:

\begin{enumerate}
  \item \textbf{Listing creation.}  The seller generates an AES-256-GCM key,
    encrypts the listing content, publishes the listing record to the DHT,
    and announces the listing key via broadcast \texttt{app\_message} to
    registered peers.

  \item \textbf{Bid submission.}  Each bidder creates a commitment
    $C = \text{SHA-256}(\mathit{bid} \| \mathit{nonce})$ with a
    cryptographically random 32-byte nonce, writes a bid record to the DHT,
    and announces the bid key to the listing's announcement registry.  The
    seller's node automatically places an auto-bid at the reserve price
    (party~0).

  \item \textbf{Deadline and party assignment.}  When the listing deadline
    expires, all nodes independently compute the same party assignment by
    sorting bidders by timestamp ascending (the seller's auto-bid is always
    earliest, becoming party~0).  This deterministic assignment requires no
    coordination.

  \item \textbf{MPC route exchange.}  Each party creates a Veilid private
    route and writes the route blob to their bid record's subkey~1.  Parties
    discover each other's routes via DHT reads with \texttt{app\_message}
    announcements as a fast path.  A readiness barrier
    (\texttt{MpcReady} via \texttt{app\_call}) ensures all routes are live
    before proceeding.

  \item \textbf{MPC execution.}  The \texttt{MpcOrchestrator} launches a
    local \texttt{mascot-party.x} (or \texttt{shamir-party.x}) process,
    feeds each party's bid value as secret input, and tunnels all inter-party
    TCP traffic through the Veilid \texttt{app\_call} tunnel.  The MPC
    program (Listing~\ref{lst:auction-mpc}) computes the maximum bid and
    reveals the winner ID and winning bid value only to party~0.

  \item \textbf{Winner verification.}  The seller sends a challenge to the
    claimed winner via \texttt{app\_call}.  The winner responds with their
    bid value and nonce; the seller verifies that
    $\text{SHA-256}(\mathit{bid} \| \mathit{nonce}) = C$.

  \item \textbf{Decryption key transfer.}  Upon successful verification, the
    seller sends the AES-256-GCM decryption key to the winner via
    \texttt{app\_call}.
\end{enumerate}


\begin{figure}[htbp]
\begin{lstlisting}[language=Python,basicstyle=\small\ttfamily,numbers=left,
  numberstyle=\tiny\color{gray},frame=single,xleftmargin=2em]
# n-Party Sealed-Bid Auction (First-Price)
# Party 0 = seller (reserve price). Parties 1..N = bidders.

n_parties = int(program.args[1])

# Each party inputs their bid as a secret integer
bids = Array(n_parties, sint)
for i in range(n_parties):
    bids[i] = sint.get_input_from(i)

# Find maximum bid and winner
max_bid = bids[0]
winner = sint(0)
for i in range(1, n_parties):
    is_higher = bids[i] > max_bid
    max_bid = is_higher.if_else(bids[i], max_bid)
    winner = is_higher.if_else(sint(i), winner)

# Reveal winner and bid ONLY to party 0 (seller)
winner_id = winner.reveal_to(0)
winning_bid = max_bid.reveal_to(0)
print_ln_to(0, 'Winner: Party %s', winner_id)
print_ln_to(0, 'Winning bid: %s', winning_bid)

# Each other party learns only won/lost
for party in range(n_parties):
    am_i_winner = (winner == sint(party)).reveal_to(party)
    print_ln_to(party, 'You won: %s', am_i_winner)
\end{lstlisting}
\caption{The MPC auction program (\texttt{auction\_n.mpc}), written in the
  MP-SPDZ DSL.  The comparison on line~16 and conditional select on lines
  17--18 are computed over secret-shared values - \(sint\), the provided primitive 
  for an (appropriately named) secret int within SPDZ. No party sees intermediate results. 
  The \texttt{reveal\_to(0)} calls on lines~21--22 ensure that
  only party~0 (the seller) learns the winner identity and winning bid.}
\label{lst:auction-mpc}
\end{figure}

\subsection{MPC Integration}
\label{sec:results-mpc}

\subsubsection{The Auction Program}

The \texttt{auction\_n.mpc} program (46 lines) implemented a first-price
sealed-bid auction for $N$ parties.  Each party inputs their bid as a secret
integer; the program computes the maximum and reveals (i)~the winner identity
and winning bid to party~0 only, and (ii)~a binary ``won: 0/1'' to each other
party.  The program is compiled once and executed with the appropriate party
count flag (\texttt{-N}). Note that for honest majority protocols executed within SPDZ, OpenSSL certificates
are required to be generated at compile time for each potential participant, and shared before the MPC reaches runtime. 
This is an advantage of using a dishonest majority protocol; no SSL secret shares prior to executing the MPC.

\subsubsection{Tunnel Architecture}

The MPC tunnel sent TCP over Veilid private routes.  Outgoing and incoming
traffic each ran in its own async task.  Sends were pipelined (up to eight
\texttt{app\_call} requests in flight at once) to minimize network latency.
The receiver reassembled the byte stream using sequence numbers and a reorder
buffer, since Veilid does not guarantee ordered delivery. Veilid has a limit of 32KB payloads
per communication envelope, and so the app\_calls were maxsized at 32KB each.

Table~\ref{tab:tunnel-volumes} shows the data volumes per protocol and party
count.  MASCOT's pairwise OT produces two to three orders of magnitude more
tunnel traffic than Shamir, meaning more \texttt{app\_call} round trips and
longer wall-clock times.

\begin{table}[H]
\centering
\small
\begin{tabular}{lrrrr}
\toprule
\textbf{Protocol} & $N$ & \textbf{Per-party (MB)} & \textbf{Total (MB)} &
  \textbf{App\_calls} \\
\midrule
MASCOT & $3$  & 37   & 107   & 3,416 \\
MASCOT & $6$  & 209  & 1,192 & 38,165 \\
MASCOT & $10$ & 653  & 6,213 & 198,747 \\
\midrule
Shamir & $3$  & 0.02  & 0.11  & 4 \\
Shamir & $6$  & 0.21  & 1.2   & 39 \\
Shamir & $10$ & 0.61  & 4.2   & 134 \\
Shamir & $20$ & 2.5   & 25.6  & 842 \\
\bottomrule
\end{tabular}
\caption{Data volumes tunnelled through Veilid per auction.  Each
  \texttt{app\_call} carries up to 32\,KB of payload.  MASCOT transfers
  roughly 1,000$\times$ more data than Shamir at $N{=}3$, growing to
  1,500$\times$ at $N{=}10$ due to $O(N^2)$ pairwise OT.}
\label{tab:tunnel-volumes}
\end{table}

\subsubsection{Route Lifecycle Challenges}

Veilid private routes are the single largest source of engineering complexity
in the system.  Three problems were encountered and solved:

\begin{enumerate}
  \item \textbf{Silent death.}  Routes die without notification; the only
    signal is a timeout on the next \texttt{app\_call}.  The system implemented
    reactive route refresh: if the readiness barrier detected fewer successful
    acknowledgements than expected, the node recreated its route, republished
    to the DHT, and polled peers' DHT entries for updated blobs.

  \item \textbf{Single active route.}  Veilid supports only one active
    private route per node.  Creating a second route invalidates the first.
    This means MPC sessions must be sequential, never concurrent.  When
    multiple auctions expire simultaneously, expired listings are sorted
    deterministically by key to ensure all nodes process them in the same
    order.

  \item \textbf{Self-import requirement.}  After creating a private route
    via \texttt{new\_private\_route()}, the node must call
    \texttt{import\_remote\_private\_route()} on its own blob.  Without this
    step, only one of $N$ created routes is deliverable.  This behaviour is
    undocumented but was confirmed by the Veilid core developers.
\end{enumerate}

These workarounds add approximately 200 lines of code.  Veilid~0.6.0 will
introduce DHT-published route blobs that are automatically fetched and
imported on use, eliminating the need for manual route
management~\cite{rioux2025_dht_patterns}.

\subsection{User Interface}
\label{sec:results-ui}

The Dioxus desktop application provided four main views: a marketplace
browser showing active listings, a listing creation form with deadline and
reserve price fields, a bid placement dialog, and an auction results view
showing the winner notification or decryption key.  The UI communicated with
the auction coordination layer through shared \texttt{Arc<Mutex<AppState>{}>}
state, polled by the Dioxus render loop.

The UI was not a focus of the project; it served as a thin presentation layer
demonstrating that the underlying protocol is usable interactively.  A
separate headless binary (\texttt{market-headless}) provided the same auction
logic without the UI, enabling automated E2E testing and benchmarking.  The
headless binary accepts listing parameters, bid values, and deadlines via
command-line arguments and environment variables, writes structured output
(JSON events) to stdout, and exits after the auction completes or times out.

\subsection{Testing and Verification}
\label{sec:results-testing}

\subsubsection{Test Coverage}

Table~\ref{tab:test-matrix} summarises the test suite.

\begin{table}[H]
\centering
\small
\begin{tabular}{llrl}
\toprule
\textbf{Level} & \textbf{Scope} & \textbf{Count} & \textbf{Runtime} \\
\midrule
Unit        & Modules (traits, serialisation, crypto) & 169 & $< 1$\,s \\
Integration & Full $N$-party auctions (mock DHT)      & 63  & $< 1$\,s \\
E2E         & Playground devnet, real Veilid networking & 3           & 3--8\,min each \\
\bottomrule
\end{tabular}
\caption{Test matrix.}
\label{tab:test-matrix}
\end{table}

\subsubsection{Integration Test Results}

The mock-based integration tests exercise the auction state machine with
3, 5, and 10 parties, covering:
\begin{itemize}[noitemsep]
  \item Highest bid wins and receives decryption key.
  \item Tied bids resolved by earliest timestamp.
  \item Reserve price enforced (auto-bid wins if no bidder exceeds it).
  \item Expired listings detected and processed.
  \item Commitment verification rejects incorrect nonces.
\end{itemize}

All 232 tests (unit + integration) pass consistently with zero flakiness.

\subsubsection{End-to-End Test Results}

The three E2E tests ran against a 20-node playground devnet:

\begin{itemize}[noitemsep]
  \item \textbf{Happy path} (169\,s): single 3-party auction
    completes successfully, winner receives decryption key.
  \item \textbf{Sequential} (380\,s): two auctions in sequence;
    tests route refresh between sessions.
  \item \textbf{Concurrent} (284\,s): two auctions with overlapping
    deadlines; verifies deterministic ordering prevents deadlock.
\end{itemize}

All three tests pass reliably on a clean devnet.  The devnet must be
restarted between test runs because Veilid's routing tables retain stale
entries for nodes that no longer exist; reconnecting on the same IP triggers
a one-hour penalty before the node is accepted again.  Additionally,
reconnecting with a different node identity on the same IP forces peers to
flush their cached knowledge of the previous node, and Veilid nodes
generally reject multiple node identities from a single IP address.

\subsection{Performance Benchmarks}
\label{sec:results-benchmarks}

The benchmark harness ran each configuration 5 times across three devnet
sizes (40, 60, 80 nodes), with party counts from 3 to 20.  Benchmarks were
run across two machines and results were combined:

\begin{itemize}[noitemsep]
  \item \textbf{Laptop}: AMD Ryzen AI 5 340 (12 cores) @ 4.90\,GHz,
    32\,GB DDR5, Arch Linux.
  \item \textbf{Desktop}: AMD Ryzen 7 5800X (16 threads) @ 4.85\,GHz,
    32\,GB DDR4, Arch Linux.
\end{itemize}

Figures in this section are generated from 354 trial rows
(153 direct TCP, 201 over Veilid).

% Benchmark figures for Chapter 5 — Performance Results
% Generated from bench-results/ CSVs via scripts/preprocess_benchmarks.py
% Re-run the script after adding new data, then recompile.

% Colours defined in Dissertation.tex preamble (plotblue, plotred, etc.)

\pgfplotsset{
  bench base/.style={
    width=0.88\columnwidth,
    height=6.5cm,
    grid=major,
    grid style={gray!30},
    tick label style={font=\small},
    label style={font=\small},
    legend style={font=\footnotesize, draw=none, fill=white, fill opacity=0.8},
    every axis plot/.append style={thick, mark size=2pt},
  },
}

%% ----------------------------------------------------------------
%% Figure 1: Direct MPC baseline (no Veilid)
%% ----------------------------------------------------------------
\begin{figure}[htbp]
  \centering
  \begin{tikzpicture}
    \begin{axis}[
      bench base,
      xlabel={Number of parties ($N$)},
      ylabel={Wall-clock time (s)},
      ymode=log,
      log basis y=10,
      xtick={3,4,5,6,8,10,15,20},
      legend pos=north west,
    ]
    \addplot+[color=plotred, mark=square*,
      error bars/.cd, y dir=both, y explicit]
      table[x=parties, y=wall_median, y error=wall_stddev]
      {Figures/data/direct_mascot.dat};
    \addlegendentry{MASCOT}

    \addplot+[color=plotblue, mark=triangle*,
      error bars/.cd, y dir=both, y explicit]
      table[x=parties, y=wall_median, y error=wall_stddev]
      {Figures/data/direct_shamir.dat};
    \addlegendentry{Shamir}

    \addplot+[color=plotgrey, mark=diamond*, only marks, mark size=3pt]
      table[x=parties, y=wall_median]
      {Figures/data/direct_rep-ring.dat};
    \addlegendentry{Rep-ring ($N{=}3$ only)}
    \end{axis}
  \end{tikzpicture}
  \caption{Direct MPC computation time (localhost TCP, no Veilid overlay).
    MASCOT scales quadratically with party count due to pairwise OT setup;
    Shamir scales linearly.  Error bars show $\pm 1$ standard deviation
    over 10 iterations.}
  \label{fig:direct-mpc-baseline}
\end{figure}

%% ----------------------------------------------------------------
%% Figure 2: Veilid auction total time — MASCOT
%% ----------------------------------------------------------------
\begin{figure}[htbp]
  \centering
  \begin{tikzpicture}
    \begin{axis}[
      bench base,
      xlabel={Number of parties ($N$)},
      ylabel={Total auction time (s)},
      xtick={3,4,5,6,8,10,12},
      legend pos=north west,
      ymin=0,
    ]
    \addplot+[color=plotblue, mark=*,
      error bars/.cd, y dir=both, y explicit]
      table[x=parties, y=total_median, y error=total_stddev,
            restrict expr to domain={\thisrow{n_success}}{1:99}]
      {Figures/data/veilid_mascot_40.dat};
    \addlegendentry{40-node devnet}

    \addplot+[color=plotorange, mark=square*,
      error bars/.cd, y dir=both, y explicit]
      table[x=parties, y=total_median, y error=total_stddev,
            restrict expr to domain={\thisrow{n_success}}{1:99}]
      {Figures/data/veilid_mascot_60.dat};
    \addlegendentry{60-node devnet}

    \addplot+[color=plotred, mark=triangle*,
      error bars/.cd, y dir=both, y explicit]
      table[x=parties, y=total_median, y error=total_stddev,
            restrict expr to domain={\thisrow{n_success}}{1:99}]
      {Figures/data/veilid_mascot_80.dat};
    \addlegendentry{80-node devnet}
    \end{axis}
  \end{tikzpicture}
  \caption{Total auction time over Veilid using MASCOT
    (median of successful runs, $\pm 1\sigma$).
    Route exchange and MPC tunnel latency dominate.  Larger devnets
    reduce success rates (see Figure~\ref{fig:success-rate}) but
    successful runs show similar wall-clock times.}
  \label{fig:veilid-mascot-total}
\end{figure}

%% ----------------------------------------------------------------
%% Figure 3: Veilid auction total time — Shamir
%% ----------------------------------------------------------------
\begin{figure}[htbp]
  \centering
  \begin{tikzpicture}
    \begin{axis}[
      bench base,
      xlabel={Number of parties ($N$)},
      ylabel={Total auction time (s)},
      xtick={3,4,5,6,8,10,15,20},
      legend pos=north west,
      ymin=0,
    ]
    \addplot+[color=plotblue, mark=*,
      error bars/.cd, y dir=both, y explicit]
      table[x=parties, y=total_median, y error=total_stddev,
            restrict expr to domain={\thisrow{n_success}}{1:99}]
      {Figures/data/veilid_shamir_40.dat};
    \addlegendentry{40-node devnet}

    \addplot+[color=plotorange, mark=square*,
      error bars/.cd, y dir=both, y explicit]
      table[x=parties, y=total_median, y error=total_stddev,
            restrict expr to domain={\thisrow{n_success}}{1:99}]
      {Figures/data/veilid_shamir_60.dat};
    \addlegendentry{60-node devnet}

    \addplot+[color=plotred, mark=triangle*,
      error bars/.cd, y dir=both, y explicit]
      table[x=parties, y=total_median, y error=total_stddev,
            restrict expr to domain={\thisrow{n_success}}{1:99}]
      {Figures/data/veilid_shamir_80.dat};
    \addlegendentry{80-node devnet}
    \end{axis}
  \end{tikzpicture}
  \caption{Total auction time over Veilid using Shamir secret sharing
    (median of successful runs, $\pm 1\sigma$).
    The 40-node devnet has insufficient relay capacity to maintain
    routes for all 20 parties ($N{=}20$: 0/10 success).  Only
    successful iterations are included; timeouts are excluded.}
  \label{fig:veilid-shamir-total}
\end{figure}

%% ----------------------------------------------------------------
%% Figure 4: Time breakdown — route exchange vs MPC (Shamir, 40-node)
%% ----------------------------------------------------------------
\begin{figure}[htbp]
  \centering
  \begin{tikzpicture}
    \begin{axis}[
      bench base,
      xlabel={Number of parties ($N$)},
      ylabel={Time (s)},
      ybar stacked,
      bar width=10pt,
      xtick={3,4,5,6,8,10,15,20},
      ymin=0,
      legend style={at={(0.02,0.98)}, anchor=north west},
    ]
    \addplot+[fill=plotblue!70, draw=plotblue,
      restrict expr to domain={\thisrow{n_success}}{1:99}]
      table[x=parties, y=route_median]
      {Figures/data/veilid_shamir_40.dat};
    \addlegendentry{Route exchange}

    \addplot+[fill=plotred!70, draw=plotred,
      restrict expr to domain={\thisrow{n_success}}{1:99}]
      table[x=parties, y=mpc_median]
      {Figures/data/veilid_shamir_40.dat};
    \addlegendentry{MPC computation}
    \end{axis}
  \end{tikzpicture}
  \caption{Time breakdown for Shamir auctions on a 40-node devnet.
    Route exchange through Veilid private routes dominates total time;
    the MPC computation itself (tunnelled over Veilid \texttt{app\_call})
    accounts for less than 10\% at most party counts.}
  \label{fig:time-breakdown}
\end{figure}


%% ----------------------------------------------------------------
%% Figure 6: Overhead factor — cost of privacy
%% ----------------------------------------------------------------
\begin{figure}[htbp]
  \centering
  \begin{tikzpicture}
    \begin{axis}[
      bench base,
      xlabel={Number of parties ($N$)},
      ylabel={Time (s)},
      ymode=log,
      log basis y=10,
      xtick={3,4,5,6,8,10,15},
      legend pos=north west,
    ]
    % Direct baselines (dashed)
    \addplot+[color=plotred, mark=square*, dashed]
      table[x=parties, y=direct_median]
      {Figures/data/overhead_mascot.dat};
    \addlegendentry{MASCOT direct}

    \addplot+[color=plotblue, mark=triangle*, dashed]
      table[x=parties, y=direct_median]
      {Figures/data/overhead_shamir.dat};
    \addlegendentry{Shamir direct}

    % Veilid MPC time (solid)
    \addplot+[color=plotred, mark=square*]
      table[x=parties, y=veilid_mpc_median]
      {Figures/data/overhead_mascot.dat};
    \addlegendentry{MASCOT over Veilid}

    \addplot+[color=plotblue, mark=triangle*]
      table[x=parties, y=veilid_mpc_median]
      {Figures/data/overhead_shamir.dat};
    \addlegendentry{Shamir over Veilid}
    \end{axis}
  \end{tikzpicture}
  \caption{MPC computation time: direct TCP vs.\ tunnelled through
    Veilid private routes (40-node devnet, median of successful runs).
    The Veilid tunnel adds approximately $15{-}19\times$ overhead for
    MASCOT and $10{-}30\times$ for Shamir, attributable to per-round
    \texttt{app\_call} latency through onion-routed hops.}
  \label{fig:overhead}
\end{figure}

%% ----------------------------------------------------------------
%% Figure 7: Bandwidth — MASCOT vs Shamir (direct)
%% ----------------------------------------------------------------
\begin{figure}[htbp]
  \centering
  \begin{tikzpicture}
    \begin{axis}[
      bench base,
      xlabel={Number of parties ($N$)},
      ylabel={Data sent per party (MB)},
      ymode=log,
      log basis y=10,
      xtick={3,4,5,6,8,10,15,20},
      legend pos=north west,
    ]
    \addplot+[color=plotred, mark=square*]
      table[x=parties, y=data_mb]
      {Figures/data/direct_mascot.dat};
    \addlegendentry{MASCOT}

    \addplot+[color=plotblue, mark=triangle*]
      table[x=parties, y=data_mb]
      {Figures/data/direct_shamir.dat};
    \addlegendentry{Shamir}
    \end{axis}
  \end{tikzpicture}
  \caption{Per-party data sent during MPC computation (direct TCP).
    MASCOT's pairwise oblivious transfer requires $O(N^2)$ bandwidth
    growth (37\,MB at $N{=}3$ to 653\,MB at $N{=}10$), while Shamir's
    secret-sharing approach grows linearly
    (0.02\,MB at $N{=}3$ to 1.4\,MB at $N{=}20$).  MASCOT benchmarks
    beyond $N{=}12$ were not feasible: the host desktop environment
    ran out of memory under the combined load of devnet nodes and
    MPC processes.}
  \label{fig:bandwidth}
\end{figure}

%% ----------------------------------------------------------------
%% Figure 8: Success rate — MASCOT vs Shamir across devnet sizes
%% ----------------------------------------------------------------
\begin{figure}[htbp]
  \centering
  \begin{tikzpicture}
    \begin{axis}[
      bench base,
      xlabel={Number of parties ($N$)},
      ylabel={Success rate},
      ymin=0, ymax=1.1,
      xtick={3,4,5,6,8,10,12,15,20},
      ytick={0,0.2,0.4,0.6,0.8,1.0},
      yticklabel={\pgfmathprintnumber\tick},
      legend pos=south west,
    ]
    \addplot+[color=plotred, mark=square*]
      table[x=parties, y=rate_40]
      {Figures/data/success_mascot.dat};
    \addlegendentry{MASCOT 40-node}

    \addplot+[color=plotblue, mark=triangle*]
      table[x=parties, y=rate_40]
      {Figures/data/success_shamir.dat};
    \addlegendentry{Shamir 40-node}

    \addplot+[color=plotred, mark=square*, dashed]
      table[x=parties, y=rate_80,
            restrict expr to domain={\thisrow{rate_80}}{0:1}]
      {Figures/data/success_mascot.dat};
    \addlegendentry{MASCOT 80-node}

    \addplot+[color=plotblue, mark=triangle*, dashed]
      table[x=parties, y=rate_80,
            restrict expr to domain={\thisrow{rate_80}}{0:1}]
      {Figures/data/success_shamir.dat};
    \addlegendentry{Shamir 80-node}
    \end{axis}
  \end{tikzpicture}
  \caption{Success rate by party count and devnet size.
    MASCOT maintained 100\% success on the 40-node devnet up to
    $N{=}12$ but degraded on larger devnets.
    Shamir scaled further in party count but dropped to 0\%
    at $N{=}20$ on 40 nodes.}
  \label{fig:success-rate}
\end{figure}

%% ----------------------------------------------------------------
%% Figure 9: MPC self-time vs wall-time (tunnel overhead scatter)
%% ----------------------------------------------------------------
\begin{figure}[htbp]
  \centering
  \begin{tikzpicture}
    \begin{axis}[
      bench base,
      xlabel={MPC self-reported time (s)},
      ylabel={MPC wall-clock time over Veilid (s)},
      xmin=0, ymin=0,
      legend pos=north west,
    ]
    % y=x reference line (no overhead)
    \addplot+[color=plotgrey, no marks, dashed, thin,
      domain=0:130] {x};
    \addlegendentry{$y = x$ (no overhead)}

    \addplot+[color=plotred, mark=square*, only marks, mark size=1.5pt,
      restrict expr to domain={\thisrow{mpc_median}}{0.01:999}]
      table[x=mpc_median, y=mpc_median,
            y expr=\thisrow{total_median}-\thisrow{route_median}]
      {Figures/data/veilid_mascot_40.dat};
    \addlegendentry{MASCOT}

    \addplot+[color=plotblue, mark=triangle*, only marks, mark size=1.5pt,
      restrict expr to domain={\thisrow{mpc_median}}{0.01:999}]
      table[x=mpc_median, y=mpc_median,
            y expr=\thisrow{total_median}-\thisrow{route_median}]
      {Figures/data/veilid_shamir_40.dat};
    \addlegendentry{Shamir}
    \end{axis}
  \end{tikzpicture}
  \caption{MPC computation time vs.\ total non-route time
    (40-node devnet).  Points above the $y{=}x$ line indicate
    overhead from Veilid tunnel latency.  MASCOT incurs larger
    absolute overhead due to its higher round count.}
  \label{fig:mpc-overhead-scatter}
\end{figure}


\subsubsection{Direct MPC Baseline}

Figure~\ref{fig:direct-mpc-baseline} shows MPC computation time on direct
TCP (localhost, no Veilid overlay).  MASCOT scales quadratically with party
count due to pairwise OT setup: 0.21\,s at $N{=}3$, rising to 5.83\,s at
$N{=}10$.  Shamir scales linearly: 0.03\,s at $N{=}3$ to 0.58\,s at
$N{=}20$.  The performance difference is roughly $30\times$ at $N{=}10$.

Figure~\ref{fig:bandwidth} shows the bandwidth cost: MASCOT sends 37\,MB per
party at $N{=}3$ and 653\,MB at $N{=}10$ ($O(N^2)$ growth), while Shamir
sends 0.02\,MB at $N{=}3$ and 2.5\,MB at $N{=}20$ (linear growth).

\subsubsection{Veilid Tunnel Overhead}

Figure~\ref{fig:overhead} compares direct TCP with Veilid-tunnelled MPC
computation time.  The Veilid tunnel added approximately $15{-}19\times$
overhead for MASCOT and $10{-}30\times$ for Shamir.  This section analyses
where that overhead comes from.

\textbf{Per-round latency.}  Each MPC communication round requires at least
one \texttt{app\_call} round trip through Veilid's onion-routed private routes.
Each private route traverses two relay hops (the default on the devnet), so a
single \texttt{app\_call} traverses at minimum four hops (sender $\to$ two
relays $\to$ recipient, then back).  At each hop, Veilid performs
Curve25519/XSalsa20-Poly1305 decryption on the envelope, routes the payload
to the next hop, and re-encrypts.  Table~\ref{tab:per-round-latency} shows
the resulting per-round latency for both protocols.

\begin{table}[H]
\centering
\small
\begin{tabular}{lrrrrr}
\toprule
\textbf{Protocol} & $N$ & \textbf{Rounds} & \textbf{Direct (ms/round)} &
  \textbf{Veilid (ms/round)} & \textbf{Slowdown} \\
\midrule
MASCOT  & $3$  & 257  & 0.65  & 17.9  & $28\times$ \\
MASCOT  & $6$  & 963  & 1.28  & 22.2  & $17\times$ \\
MASCOT  & $10$ & 2483 & 2.32  & 36.4  & $16\times$ \\
Shamir  & $3$  & 35   & 0.05  & 10.9  & $209\times$ \\
Shamir  & $10$ & 122  & 0.38  & 29.3  & $77\times$ \\
Shamir  & $20$ & 242  & 1.26  & ---   & ---  \\
\bottomrule
\end{tabular}
\caption{Per-round MPC latency: direct TCP vs.\ Veilid tunnel (40-node
  devnet).  Per-round latency is computed as MPC wall-clock time divided
  by round count.  Shamir $N{=}20$ is absent because no trial succeeded
  on the 40-node devnet.}
\label{tab:per-round-latency}
\end{table}

The per-round overhead explains why the \emph{aggregate} overhead factor
differs between MASCOT and Shamir.  MASCOT performs substantial local
computation per round (OT extension, MAC generation), so the fixed
${\thicksim}$18\,ms network latency is a smaller fraction of each round's total
time.  Shamir's local computation is sub-microsecond, so the same
${\thicksim}$18\,ms network latency accounts for nearly all of each round's
cost, producing a higher per-round slowdown factor.  In absolute terms,
Shamir remains far faster due to fewer rounds and less data per round.

\textbf{Tunnel throughput.}  The tunnel's effective throughput is limited by
\texttt{app\_call} round-trip time and pipeline depth.  With a pipeline depth
of 8 and ${\thicksim}$18\,ms per \texttt{app\_call}, the theoretical maximum
throughput is $8 \times 32\,\text{KB} / 18\,\text{ms} \approx 14\,\text{MB/s}$.
In practice, the tunnel sustained approximately 2.5\,MB/s for MASCOT $N{=}3$
(109\,MB in 44\,s of MPC time) and 3\,KB/s for Shamir $N{=}3$ (116\,KB in
0.38\,s).  The gap between theoretical and observed throughput for MASCOT
arises because not all 8 pipeline slots are always filled: the MP-SPDZ process
alternates between sending data and waiting for responses in each round.

\textbf{Cryptographic cost per hop.}  Each Veilid relay hop performs
envelope decryption (XSalsa20-Poly1305) and re-encryption for the next hop.
On the benchmark hardware (Ryzen AI 5 340), XSalsa20-Poly1305 processes
${\thicksim}$1\,GB/s, so the per-hop crypto cost for a 32\,KB \texttt{app\_call}
payload is ${\thicksim}$32\,\textmu{}s. This is negligible compared to the network
round-trip latency.  The overhead is therefore dominated by \emph{network
traversal} (hop-by-hop forwarding, connection multiplexing, congestion
backoff), not by cryptographic computation.

\subsubsection{Total Auction Time}

Figures~\ref{fig:veilid-mascot-total} and~\ref{fig:veilid-shamir-total} show
total end-to-end auction time (route exchange + MPC + verification).
Route exchange dominated total time at every party count.
Figure~\ref{fig:time-breakdown} shows that for Shamir on a 40-node devnet,
MPC computation accounted for less than 10\% of wall-clock time at all party
counts tested.  Even for MASCOT, where the tunnel carries hundreds of
megabytes, route exchange still accounted for 40--65\% of total time.

The route exchange phase consists of four sequential steps, each involving
Veilid operations:
\begin{enumerate}[noitemsep]
  \item \textbf{Route construction} (${\thicksim}$5--15\,s): each party calls
    \texttt{new\_custom\_private\_route()} and self-imports the blob.  This
    involves Diffie--Hellman key exchanges with two relay nodes.
  \item \textbf{DHT publishing} (${\thicksim}$2--5\,s): the route blob is
    written to the party's bid DHT record (subkey~1).  DHT writes
    propagate through the overlay's Kademlia routing.
  \item \textbf{Route discovery} (${\thicksim}$10--40\,s, scales with $N$):
    each party reads all other parties' DHT records to collect route blobs.
    The $N{-}1$ DHT reads are sequential because each read forces a network
    round trip to the DHT record's close nodes.
  \item \textbf{Readiness barrier} (${\thicksim}$5--30\,s): each party sends an
    \texttt{MpcReady} message via \texttt{app\_call} to every other party's
    imported route, confirming all routes are live.  Failed deliveries trigger
    a reactive refresh cycle: recreate the local route, republish to DHT,
    re-poll peers.
\end{enumerate}

Route discovery dominates because it requires $O(N)$ sequential DHT reads,
each taking 1--4\,s depending on how quickly the DHT close nodes respond.
At $N{=}10$, nine DHT reads at ${\thicksim}$3\,s each account for ${\thicksim}$27\,s
of the 213\,s median route exchange time; the remainder is consumed by the
readiness barrier and retry cycles.

This analysis suggests that reducing route exchange time (through
persistent routes, DHT caching, or parallel route discovery) would
have a larger impact on total auction time than optimising the MPC tunnel
itself.

\subsubsection{Reliability and Success Rates}

Not all trials succeeded.  Figure~\ref{fig:success-rate} shows success rates
by party count and devnet size.  On the 40-node devnet, MASCOT achieved
100\% success up to $N{=}12$ (the largest party count tested).  Shamir
maintained high success rates up to $N{=}10$ (70\%) but dropped to 0\% at
$N{=}20$.  On the 80-node devnet, both protocols degraded: MASCOT fell to
10\% at $N{=}10$ and Shamir to 40\% at $N{=}20$.

Three distinct failure modes were observed:
\begin{enumerate}[noitemsep]
  \item \textbf{Route construction timeout}: the node could not find a
    suitable relay node to form a 1-hop private route within the 50\,s
    deadline (Veilid defaults to 2 hops; this was reduced to 1 for the
    smaller devnet topology).  This occurred more frequently on larger
    devnets where routing table convergence took longer.
  \item \textbf{Dead route during MPC}: a relay node in an active private
    route became unreachable mid-computation.  Because Veilid does not
    notify applications of route death, the tunnel detected this only after
    a 10\,s \texttt{app\_call} timeout, at which point the MPC process had
    already stalled.
  \item \textbf{MPC process timeout}: the 600\,s global deadline was
    exceeded, typically because multiple dead-route recovery cycles consumed
    the time budget.
\end{enumerate}

All failures were attributable to Veilid route instability under load, not
to bugs in the auction protocol or MPC program.  The MP-SPDZ processes
themselves never produced incorrect results when they completed.

\subsubsection{Devnet Size Effects}

Larger devnets (60, 80 nodes) showed \emph{worse} performance, not better.
This is counterintuitive: one might expect more relay nodes to provide more
routing options and thus better availability.  Three factors explain the
degradation:

\begin{enumerate}[noitemsep]
  \item \textbf{Routing table convergence.}  Veilid's Kademlia-based routing
    tables grow with network size.  On a freshly started 80-node devnet, each
    node must discover and validate up to 79 peers via iterative lookups.
    During this convergence period, route construction attempts may select
    relay nodes that have not yet stabilised their own routing tables.
  \item \textbf{Resource contention.}  All devnet nodes ran on the same
    physical machine, sharing CPU, memory, and network stack.  At 80 nodes,
    each Veilid process consumed ${\thicksim}$50\,MB of RAM and maintained
    ${\thicksim}$40 UDP connections, totalling ${\thicksim}$4\,GB RAM and
    ${\thicksim}$3,200 sockets.  Context switching and socket buffer contention
    added artificial latency not present in a real distributed deployment.
  \item \textbf{Relay selection quality.}  Veilid selects relay nodes based
    on latency and reliability metrics that take time to accumulate.  On
    larger devnets, the probability of selecting an under-performing relay
    increases during the initial convergence period.
\end{enumerate}

On the public Veilid network, with thousands of nodes across diverse hardware
and with stable long-running relays, the behaviour would likely differ.
However, public network testing was not practical: Veilid enforces
one-hour bans per IP for misconfigured nodes~\cite{rioux2026_ip_punishment},
which occurs frequently during development, and running behind CGNAT meant
nodes did not meet Veilid's reachability requirements for full relay
participation, motivating the devnet approach.

\subsubsection{Hardware and Software Limits}

The practical limit for MASCOT over Veilid on this devnet is approximately
$N{=}6$ (100\% success rate on the 40-node devnet, degrading on larger
devnets).  Beyond this, the $O(N^2)$ bandwidth overwhelms the tunnel: at
$N{=}10$, MASCOT sends ${\thicksim}$637\,MB per party through the tunnel,
requiring ${\thicksim}$2,500 \texttt{app\_call} round trips.

Shamir scales further because its linear bandwidth means the tunnel carries
${\thicksim}$2.3\,MB rather than 637\,MB at comparable party counts.  Shamir
$N{=}15$ achieved 70\% success on the 40-node devnet; $N{=}20$ succeeded on
60 and 80-node devnets (30--40\%) but not on 40 nodes, likely due to
insufficient relay diversity at the smaller size.

The bottleneck is not CPU or memory but \emph{Veilid route latency}.
Table~\ref{tab:warm-cold} illustrates the difference between warm and cold
devnets.

\begin{table}[H]
\centering
\small
\begin{tabular}{lrrrr}
\toprule
\textbf{Condition} & \textbf{MPC time} & \textbf{Route time} &
  \textbf{Total} & \textbf{ms/round} \\
\midrule
40-node devnet (median) & 4.61\,s  & 55.66\,s & 99.97\,s & 17.9 \\
Direct TCP (no Veilid)  & 0.20\,s  & ---      & 0.20\,s  & 0.78 \\
\bottomrule
\end{tabular}
\caption{MASCOT $N{=}3$: Veilid tunnel vs.\ direct TCP (40-node devnet,
  10 successful trials).  The tunnel adds 17.1\,ms/round of overhead,
  dominated by Veilid relay latency.}
\label{tab:warm-cold}
\end{table}

The Veilid tunnel adds ${\thicksim}$18\,ms of network round-trip time per
\texttt{app\_call}, compared to sub-millisecond on localhost TCP.

\subsection{Security Analysis}
\label{sec:results-security}

This section evaluates the system against the security policy from
Section~\ref{sec:intro-security-policy}, identifying what is protected,
what leaks, and the severity of each residual channel.

\subsubsection{Threat Model}

Four adversary types are considered:
\begin{itemize}[noitemsep]
  \item \textbf{Honest-but-curious bidder}: follows the protocol but attempts
    to learn other bids from observable data.
  \item \textbf{Malicious seller}: controls the listing, decryption key, and
    MPC party~0 output.
  \item \textbf{Passive network observer}: monitors DHT writes and Veilid
    message patterns without participating.
  \item \textbf{Colluding parties}: multiple participants share information to
    break the MPC threshold.
\end{itemize}

\subsubsection{Timing Side Channel on Party Assignment}

Party order is determined by bid timestamp.  A network observer monitoring DHT
writes can observe bid timing, inferring party assignments.  If the adversary
knows the mapping between party IDs and bidders, they can correlate MPC
outputs with specific participants.

\textbf{Severity}: Medium.  \textbf{Mitigation}: None currently.  Possible
fixes include random party assignment (requiring a pre-committed randomness
source) or padding bid submissions to fixed intervals.

\subsubsection{Seller as Single Point of Trust}

The seller (party~0) sees the winner identity and winning bid value, controls
decryption key release, and can refuse verification (griefing) or collude with
a preferred bidder.  This is the most significant trust assumption in the
system.

Compared to the Danish sugar-beet auction~\cite{bogetoft2009_mpc_live}, where
the auctioneer was a trusted cooperative, this system gives the seller
\emph{less} power (they see only the winning bid, not all bids) but
\emph{more} ability to grief (they can withhold the key).

\textbf{Severity}: High (by design).  \textbf{Mitigation}: Threshold
decryption (key shared among MPC parties) would remove this assumption;
see Section~\ref{sec:conc-future}.

\subsubsection{Winner Identity Linkage}

The challenge-response protocol reveals the winner to the seller before the
decryption key is sent.  This is necessary for protocol correctness: the
seller must verify the bid commitment before releasing the key.

\textbf{Severity}: Medium (by design).  The PDD proposed 1-out-of-$N$
oblivious transfer to avoid this linkage; however, since the seller already
learns the winner as MPC party~0, OT would not provide additional privacy.

\subsubsection{DHT Metadata Leakage}

Bid record keys are publicly announced via the listing's announcement
registry.  The \emph{number} and \emph{timing} of bids are observable by any
DHT reader, even though bid \emph{values} are hidden behind commitments.
This is a form of participation metadata leakage.

\textbf{Severity}: Low.  \textbf{Partial mitigation}: Veilid's private
routing reduces IP-level correlation between a bid record write and a
specific network identity.

\subsubsection{MPC Route Correlation}

MPC route blobs are written to bid record subkeys before MPC execution.
A network observer could correlate route creation timing with auction
participation, partially deanonymising MPC parties at the network layer.

\textbf{Severity}: Medium.  \textbf{Partial mitigation}: Veilid's onion
routing limits direct observation (default hop count = 2; benchmarks used
hop count = 1 for performance).  However, the anonymity set is small on a
devnet.

\subsubsection{Small-$N$ Security Considerations}

The choice between MASCOT and Shamir has direct security implications at
small party counts:

\begin{itemize}[noitemsep]
  \item \textbf{MASCOT} ($N{=}3$): tolerates up to 2 corrupt parties.  Even
    with 2 colluding bidders, the protocol remains secure,only the
    colluding parties' own inputs are ``lost'' (they already know their
    own bids).
  \item \textbf{Shamir} ($N{=}3$): requires $< 2$ corrupt parties (i.e.,
    at most 1).  A single colluding bidder plus the seller can
    reconstruct all shares, completely breaking privacy.
\end{itemize}

For small auctions (the common case in a marketplace), MASCOT's stronger
security guarantee is significant, despite its higher computational cost.

\subsubsection{Summary of Side Channels}

Table~\ref{tab:side-channels} summarises the identified information leakage
channels.

\begin{table}[H]
\centering
\small
\begin{tabular}{p{3.8cm}ccp{3.5cm}}
\toprule
\textbf{Side channel} & \textbf{Severity} & \textbf{Mitigated?} & \textbf{Exploitable by} \\
\midrule
Bid timing $\rightarrow$ party ID     & Medium & No        & Network observer \\
Seller sees winner + bid               & High   & By design & Malicious seller \\
Winner identity linkage                & Medium & By design & Seller \\
DHT participation metadata             & Low    & Partially & Any DHT reader \\
MPC route correlation                  & Medium & Partially & Network observer \\
Shamir threshold ($N{=}3$)             & High   & Use MASCOT & Colluding parties \\
\bottomrule
\end{tabular}
\caption{Summary of residual information leakage channels.}
\label{tab:side-channels}
\end{table}
